\section{Introduction}

The Food and Drug Administration is the gatekeeper of the U.S. pharmaceutical market: its approval decisions determine which drugs may be sold, when they reach patients, and what evidence of safety and efficacy is required before market entry \parencite{carpenterINTRODUCTIONGatekeeper2010}. Since 1992, the agency has been partly funded by industry user fees collected in exchange for binding review-time targets. The Prescription Drug User Fee Act (PDUFA) of 1992 and the Generic Drug User Fee Amendments (GDUFA) of 2012 are widely credited with substantially accelerating drug review and changing the incentive environment facing both the FDA and the firms it regulates \parencite{berndtIndustryFundingFDA2005, olsonPDUFAInitialUS2009, berndtGenericDrugUser2018}. They also coincide, in the United States, with one of the largest public-health crises of the past three decades: drug overdose deaths reached approximately 107{,}000 in 2021, with prescription controlled substances playing a prominent role throughout the trajectory of the epidemic \parencite{markDrugOverdoseDeaths2024}.

A large literature studies the consequences of regulatory acceleration. \textcite{berndtIndustryFundingFDA2005} document that PDUFA substantially reduced NDA review times and increased approval rates. \textcite{carpenterDrugreviewDeadlinesSafety2008} show that drugs approved just before PDUFA deadlines exhibit higher rates of subsequent safety-related labeling actions than those approved just after. \textcite{philipsonCostbenefitAnalysisFDA2008} build a welfare framework around the speed--safety tradeoff. \textcite{olsonPDUFAInitialUS2009} and \textcite{olsonFirmCharacteristicsSpeed1997} examine launch behavior and firm-side heterogeneity in pre-PDUFA review times. The accelerated-approval and coverage literature, including \textcite{darrowFDAApprovalRegulation2020} and \textcite{gelladAcceleratedApprovalExpensive2017}, documents how shorter review reduces clinical evidence requirements at approval. A separate body of work in the economics of controlled substances studies the supply side of misuse and diversion: \textcite{alpertSupplySideDrugPolicy2018a}, for example, show that the 2010 abuse-deterrent reformulation of OxyContin reduced prescription opioid misuse but shifted users into illicit markets.

These two literatures have not been joined. Existing studies of user-fee reform identify effects on review times, launch timing, post-market safety, and welfare, but do not ask whether either reform changed the \textit{composition} of approved drugs along the controlled-substance margin. The omission matters because the standard economic argument for why user-fee acceleration should reshape approval composition is straightforward. Faster review compresses the window between approval and patent expiration and therefore raises the present value of approval; the value of that compression is greatest for drugs with the largest expected revenue streams. Branded extended-release and abuse-deterrent controlled substances --- particularly Schedule~II opioid formulations --- occupy precisely this segment of the market: high list prices, long marketing horizons, and limited generic competition during patent protection \parencite{powellHowIncreasingMedical2020}. If user-fee acceleration tilts the field toward high-margin pipelines, controlled substances should be among the most affected categories. The same incentive logic applies, with sign reversed, to the generic-drug margin under GDUFA, where Schedule~II opioid generics historically face higher manufacturing and compliance costs than non-controlled generics. The null hypothesis is straightforward, theory-relevant, and economically interesting: the user-fee channel could plausibly have shifted the composition of approvals.

This paper provides the first systematic test of that hypothesis. Using the complete Drugs@FDA database linked to DEA scheduling records (1939--2025, $N = 25{,}908$ original approvals), I estimate interrupted time series specifications, Poisson rate models with total approvals as exposure, and stacked difference-in-differences specifications around both PDUFA (1992) and GDUFA (2012). Each specification is run on the application-type subsample where the corresponding policy applies most directly --- new drug applications (NDAs) for PDUFA, abbreviated new drug applications (ANDAs) for GDUFA --- and is supplemented with rate-versus-count diagnostics, group-specific linear time trends, and event-study DD designs.

The headline finding is null. In the most flexible specifications --- those that allow controlled-substance and non-controlled-substance approval streams to follow their own linear trends --- the estimated effect of either reform on the CS share of approvals is indistinguishable from zero, with confidence intervals that span both small and economically meaningful effects on either side of the null. For PDUFA on the NDA pathway, the post-policy incidence-rate ratio is $1.19$ with a 95\% confidence interval of $[0.59, 2.39]$. For GDUFA on the ANDA pathway, the corresponding IRR is $1.38$ with a 95\% confidence interval of $[0.86, 2.21]$. The point estimates flip sign across less flexible and more flexible specifications, which is itself diagnostic: when the only thing absorbing pre/post differences is a year fixed effect, post-PDUFA NDA rates appear elevated and post-GDUFA ANDA rates appear depressed, but both shifts are absorbed once group-specific trends accommodate the slower-moving secular changes in pharmaceutical innovation.

A drug-level investigation of the CS-NDA series locates the most economically interpretable confound. Forty-three CS NDAs were approved between 2009 and 2013 --- roughly twice the typical pace of the preceding decade --- and the cluster maps almost one-for-one onto the well-documented wave of branded extended-release and abuse-deterrent opioid formulations of the late 2000s: the reformulated OxyContin (2010), Exalgo (2010), Embeda (2009), Suboxone (2010), Zohydro ER (2012), and related products. The timing is roughly seventeen to twenty-one years after PDUFA enactment, which is inconsistent with a delayed regulatory-incentive response and consistent with a pharmaceutical product cycle that the FDA accommodated rather than caused. Once this episode is allowed to load on its own group-specific trend rather than on the post-PDUFA dummy, the user-fee effect collapses.

Two broader contributions follow. The first is empirical: across a comprehensive panel of FDA approvals over more than eight decades, the controlled-substance composition of the legal pharmaceutical supply does not appear to respond to user-fee-funded review acceleration. The second is interpretive. Negative results in policy economics are sometimes read as evidence that the question was poorly posed; here, the structure of the test --- a theory-relevant prediction, a comprehensive panel, multiple specifications consistent in direction --- suggests instead that the relevant margin operates elsewhere. The pharmaceutical-innovation channel that produced the late-2000s opioid wave appears to have run on firm-level investment decisions and product-cycle dynamics, with FDA review serving as a downstream accommodation rather than an upstream driver. For the design of future regulatory reforms, this implies that user-fee acceleration is unlikely, on its own, to be a useful margin for influencing the controlled-substance content of the legal pharmaceutical supply.

\paragraph{Roadmap.} Section~\ref{sec:background} reviews the institutional background, the relevant literature, and the theory of user-fee incentives. Section~\ref{sec:data} describes the data and variable construction. Section~\ref{sec:strategy} sets out the empirical strategy. Section~\ref{sec:results} reports the main and supplementary results. Section~\ref{sec:discussion} interprets the findings and discusses their limitations. Section~\ref{sec:conclusion} concludes.
